\section{12 Embedded Machine Learning}

\subsection{AI / ML Foundations}

\mult{2}

\begin{concept}{AI $\supset$ ML $\supset$ DL $\supset$ GenAI}
    \begin{itemize}
        \item \important{AI}: emulate human behaviour (learn, decide, recognise)
        \item \important{ML}: algorithms that detect patterns in data (supervised /
              unsupervised)
        \item \important{DL}: ML using multi-layer \important{neural networks}
        \item \important{GenAI}: DL that \emph{generates} content (text/images/code)
    \end{itemize}
    ML types: \important{task-driven} (supervised, predict next value),
    \important{data-driven} (unsupervised, find clusters), \important{reinforcement}
    (learn from mistakes).
\end{concept}

\begin{concept}{Feed-Forward Neural Net}
    Built from \important{neurons} in input / hidden / output layers, \important{fully
    connected}.
    \begin{itemize}
        \item knowledge stored in \important{weights $w$} and \important{bias $b$}
        \item can approximate \important{any continuous function}
        \item non-linearity from an \important{activation function}
    \end{itemize}
\end{concept}

\multend

\begin{formula}{Neuron Output \& Parameter Count}
    One neuron:
    $$a = \sigma\!\left(\sum_i w_i a_i + b\right) = \sigma(w_1 a_1 + \dots + w_n a_n + b)$$
    The bias shifts the activation threshold. MNIST net (784$\to$16$\to$16$\to$10):
    $$\underbrace{784\cdot16 + 16\cdot16 + 16\cdot10}_{\text{weights}} +
    \underbrace{16+16+10}_{\text{biases}} = \important{13\,002}\ \text{parameters}$$
    \important{Learning = finding the right weights \& biases.}
\end{formula}

\subsection{Activation Functions}

\begin{formula}{Activation Functions}
    \begin{itemize}
        \item \important{Sigmoid}: $\sigma(z)=\dfrac{1}{1+e^{-z}}$ -- "old school", not
              for hidden layers
        \item \important{Tanh}: $\dfrac{e^{z}-e^{-z}}{e^{z}+e^{-z}}$ -- small gradient
        \item \important{ReLU}: $\max(0,z)$ -- near-linear, easy to optimise, default choice
        \item \important{Leaky ReLU}: $\max(\alpha z, z)$ -- small $\alpha$ avoids ReLU's
              zero-gradient for $z\leq0$
        \item \important{Softmax} (output layer): values 0--1 summing to 1 $\to$
              \important{probabilities}
    \end{itemize}
\end{formula}

\subsection{Cost \& Training}

\mult{2}

\begin{concept}{Cost \& Gradient Descent}
    \important{Cost} = how wrong the output is vs the label; \important{lower = better}.
    Training \important{minimises} the cost by walking downhill along its gradient.
    Risk: getting stuck in \important{local minima / saddle points}.
\end{concept}

\begin{formula}{Cost / GD / Cross-Entropy}
    Sum of squared differences:
    $$C = \sum_k (\text{out}_k - \text{label}_k)^2$$
    Gradient descent step ($\alpha$ = learning rate):
    $$x \leftarrow x - \alpha\,\nabla f(x)$$
    Cross-entropy loss ($p$ true, $q$ predicted):
    $$H(p,q) = -\sum_{x} p(x)\log q(x)$$
\end{formula}

\multend

\mult{2}

\begin{definition}{Optimizers}
    \begin{itemize}
        \item \important{Gradient Descent}: basic downhill step
        \item \important{Momentum}: adds a velocity term (smooths, escapes small dips)
        \item \important{RMSProp}: per-parameter adaptive step size
        \item \important{Adam} (Adaptive Moment Est.): \important{RMSProp + Momentum},
              the usual default
    \end{itemize}
\end{definition}

\begin{definition}{Epoch / Batch / Iteration}
    \begin{itemize}
        \item \important{Epoch}: one full pass of the \emph{entire} dataset (fwd + back)
        \item \important{Batch size}: training examples per batch
        \item \important{Iterations}: batches needed for one epoch
    \end{itemize}
    $$\text{1 Epoch} = \text{Iterations} \times \text{Batch Size}$$
\end{definition}

\multend

\subsection{Keras Workflow}

\begin{definition}{Frameworks}
    \important{Keras} -- high-level API (runs on TensorFlow, \texttt{tf.keras});
    \important{TensorFlow} (Google) -- dataflow/symbolic math;
    \important{PyTorch} (Meta) -- Torch-based. All written in Python.
\end{definition}

\begin{KR}{Build / Compile / Train a Model}
\begin{lstlisting}[language=Python, style=basesmol, aboveskip=2pt, belowskip=2pt]
from tensorflow import keras
from tensorflow.keras import layers
model = keras.Sequential([                       # define
keras.Input(shape=input_shape),
layers.Flatten(),                            # 28x28 -> 784
layers.Dense(16, activation="relu", name="Layer1"),
layers.Dense(16, activation="relu", name="Layer2"),
layers.Dense(num_classes, activation="softmax"), # 10 outputs
])
model.summary()
adam = keras.optimizers.Adam(learning_rate=0.005,        # compile
beta_1=0.9, beta_2=0.999, amsgrad=False)
model.compile(optimizer=adam,
loss='sparse_categorical_crossentropy', metrics=['accuracy'])
model.fit(x_train, y_train, batch_size=50, epochs=5, verbose=1)  # train
\end{lstlisting}
    Other losses: \texttt{mse}, \texttt{mae}, \texttt{categorical\_crossentropy}.
\end{KR}

\begin{KR}{Load \& Preprocess MNIST}
\begin{lstlisting}[language=Python, style=basesmol, aboveskip=2pt, belowskip=2pt]
(x_train, y_train), (x_test, y_test) = keras.datasets.mnist.load_data()
x_train = x_train.astype("float32") / 255    # scale pixels to [0,1]
x_test  = x_test.astype("float32")  / 255
x_train = np.expand_dims(x_train, -1)        # (28,28) -> (28,28,1)
x_test  = np.expand_dims(x_test, -1)
\end{lstlisting}
    60\,000 training + 10\,000 test images.
\end{KR}

\subsection{Deployment (TensorFlow Lite)}

\begin{KR}{Convert to TFLite \& Run Inference}
    \textbf{Convert} (on host):
\begin{lstlisting}[language=Python, style=basesmol, aboveskip=2pt, belowskip=2pt]
converter = tf.lite.TFLiteConverter.from_keras_model(model)
open("model.tflite", "wb").write(converter.convert())
\end{lstlisting}
    \textbf{Inference} (on target): input = 28$\times$28 floats (0..1), output = 10
    probabilities; pick the \texttt{argmax}.
\begin{lstlisting}[language=Python, style=basesmol, aboveskip=2pt, belowskip=2pt]
interpreter = Interpreter("model.tflite")
interpreter.allocate_tensors()
in_d  = interpreter.get_input_details()[0]
out_d = interpreter.get_output_details()[0]
def predict(image):
inp = np.array(image, dtype="f").reshape([1, 28, 28, 1])
interpreter.set_tensor(in_d["index"], inp)
interpreter.invoke()
out = interpreter.get_tensor(out_d["index"])[0]
return int(np.argmax(out))           # most probable digit
\end{lstlisting}
    In C: \texttt{TfLiteModelCreateFromFile} $\to$ \texttt{...InterpreterCreate} $\to$
    \texttt{AllocateTensors} $\to$ copy input $\to$ \texttt{...Invoke} $\to$ read output.
\end{KR}

\subsection{Convolutional Neural Networks}

%% >>> IMAGE (slide 52): CNN for MNIST -- INPUT 28x28x1 -> Conv 5x5 -> MaxPool 2x2 -> Conv 5x5 -> MaxPool 2x2 -> Flatten -> FC -> 10 outputs <<<

\mult{2}

\begin{concept}{CNN}
    \begin{itemize}
        \item \important{Convolution}: filter input with a kernel $\to$
              \important{feature extraction} + data reduction
        \item \important{Pooling} (sub-sampling): max- or average-pool reduces
              dimensions \& parameters $\to$ faster training, controls
              \important{over-fitting}
    \end{itemize}
\end{concept}

\begin{formula}{Conv / Pool Output Sizes}
    Valid 5$\times$5 conv: $n \to n-4$; 2$\times$2 max-pool (stride 2): $n \to n/2$.
    \begin{itemize}
        \item 28 $\to$ \texttt{Conv5} 24 $\to$ \texttt{Pool2} 12
        \item 12 $\to$ \texttt{Conv5} 8 $\to$ \texttt{Pool2} 4
        \item flatten $4\cdot4\cdot32 = 512 \to$ Dense 10
    \end{itemize}
\end{formula}

\multend

\begin{KR}{Build a CNN in Keras}
\begin{lstlisting}[language=Python, style=basesmol, aboveskip=2pt, belowskip=2pt]
model = keras.Sequential([
keras.Input(shape=input_shape),
layers.Conv2D(16, kernel_size=(5,5), activation="relu"),
layers.MaxPooling2D(pool_size=(2,2)),
layers.Conv2D(32, kernel_size=(5,5), activation="relu"),
layers.MaxPooling2D(pool_size=(2,2)),
layers.Flatten(),
layers.Dense(num_classes, activation="softmax"),
])
# ~18'378 params, ~98.67% MNIST accuracy (vs dense-only net)
\end{lstlisting}
\end{KR}

\subsection{Image Acquisition for ML (Lab 7)}

\begin{KR}{Webcam $\to$ 28$\times$28 Pipeline}
    Host trains $\to$ \texttt{model.tflite}; target grabs \& preprocesses an image:
\begin{lstlisting}[language=bash, style=basesmol, aboveskip=2pt, belowskip=2pt]
gst-launch-1.0 v4l2src device=/dev/video0 \
  ! video/x-raw,width=640,height=480 \
  ! videorate ! "video/x-raw,framerate=1/1" \    # 1 fps
  ! videobalance contrast=2.0 \                  # increase contrast
  ! videoconvert ! "video/x-raw,format=GRAY8" \  # greyscale
  ! videocrop top=0 left=80 right=80 bottom=0 \  # 640x480 -> 480x480
  ! videoscale ! "video/x-raw,height=28,width=28" \  # -> 28x28
  ! multifilesink location="image.raw"
\end{lstlisting}
    Use \texttt{tee} to branch the stream (e.g. save both the 28$\times$28 raw and a
    JPEG preview).
\end{KR}

\subsection{Beyond: GAN, Transformers, LLMs}

\mult{2}

\begin{definition}{GAN}
    \important{Generative Adversarial Network}: a \important{generator} creates fake
    samples to fool a \important{discriminator}, which classifies real vs fake; both
    improve via back-propagation.
\end{definition}

\begin{definition}{AI Evolution}
    \important{FCN} (math functions) $\to$ \important{CNN} (2-D spatial features) $\to$
    \important{Transformers} (sequence/context) $\to$ \important{LLMs} (reasoning at
    extreme scale).
\end{definition}

\multend

\begin{definition}{LLMs (Transformer Workflow)}
    \begin{enumerate}
        \item \important{Tokenization \& embedding}: text $\to$ tokens $\to$
              high-dimensional vectors
        \item \important{Self-attention (QKV)}: Query/Key/Value vectors weigh how
              relevant surrounding tokens are
        \item \important{Probability \& softmax}: rank the vocabulary, pick the
              highest-probability next token
    \end{enumerate}
    LLMs = \important{next-token prediction}; a text subset of GenAI, trained with
    billions of parameters.
\end{definition}

\raggedcolumns
